\documentclass[11pt,a4paper]{article}

\usepackage[utf8]{inputenc}
\usepackage[T1]{fontenc}
\usepackage{amsmath,amssymb}
\usepackage{geometry}
\geometry{margin=2.5cm}

\title{Bridging the Heisenberg Spin Model to Spatial IETS $dI/dV$ Maps\\
\large Triangulene--Phenalenyl Nanographene Polymer}
\author{}
\date{\today}

\newcommand{\ket}[1]{\left|#1\right\rangle}
\newcommand{\bra}[1]{\left\langle#1\right|}
\newcommand{\braket}[2]{\left\langle#1\middle|#2\right\rangle}
\newcommand{\melem}[3]{\left\langle#1\middle|#2\middle|#3\right\rangle}
\renewcommand{\vec}[1]{\mathbf{#1}}

\begin{document}
\maketitle

\section{Introduction: from a discrete spin lattice to a spatial STM observable}

The low-energy magnetism of the phenalenyl--triangulene copolymer
$P\!-\!T\!-\!P\!-\![P\!-\!T\!-\!P]_{N-1}$ is captured by a Heisenberg Hamiltonian on
a chain of $4N$ spin-$\tfrac12$ sites. With the ordering $[P,T_a,T_b,P]$ inside
each monomer,
\begin{equation}
  H \;=\;
  \sum_{m=0}^{N-1}\!\Big[
      J_{TP}\big(\vec S_{4m}\!\cdot\!\vec S_{4m+1}
               +\vec S_{4m+2}\!\cdot\!\vec S_{4m+3}\big)
    + J_{T}\,\vec S_{4m+1}\!\cdot\!\vec S_{4m+2}\Big]
  + J_{PP}\!\sum_{m=0}^{N-2}\! \vec S_{4m+3}\!\cdot\!\vec S_{4m+4},
  \label{eq:H}
\end{equation}
with $J_{PP}=38$~meV, $J_{TP}=40$~meV and $J_{T}=-500$~meV. Exact
diagonalization yields a total-singlet ground state $\ket{0}$ ($S=0$) and, as the
first magnetic excitation, a triplet multiplet $\{\ket{1_m}\}_{m=-1,0,+1}$ ($S=1$)
separated by the spin gap $\Delta=E_1-E_0$.

Inelastic electron tunneling spectroscopy (IETS) in a low-temperature STM probes
exactly this excitation: when the bias energy $e|V|$ exceeds $\Delta$, a tunneling
electron can flip the local spin and open an inelastic channel, producing a
symmetric step in $dI/dV$ at $e|V|=\Delta$. Crucially, the \emph{intensity} of
that step is spatially resolved, because it is governed by the local spin
matrix element between $\ket{0}$ and $\ket{1_m}$ at the position of the tip. The
purpose of this note is to make explicit the bridge between the discrete spin
model~\eqref{eq:H} and the continuous $dI/dV(x,y)$ image recorded in experiment.

\section{Transition weights from the spin operators}

To second order in the tip--sample coupling, the inelastic spin-flip conductance
at site $i$ is proportional to the Fermi-golden-rule rate summed over the
degenerate final states of the excitation. For a spin-$\tfrac12$ scatterer the
relevant tunneling vertex is the local spin operator
$\vec S_i=(S_i^x,S_i^y,S_i^z)$, so the \emph{per-site transition weight} reads
\begin{equation}
  W_i \;=\; \sum_{m}\sum_{\alpha=x,y,z}
            \big|\melem{1_m}{S_i^{\alpha}}{0}\big|^{2},
  \label{eq:Wi}
\end{equation}
where the inner sum runs over the three Cartesian spin components and the outer
sum over the members $m$ of the final triplet multiplet. Writing out the three
terms for a single final state,
\begin{equation}
  W_i \;=\; \big|\melem{1}{S_i^{x}}{0}\big|^{2}
          + \big|\melem{1}{S_i^{y}}{0}\big|^{2}
          + \big|\melem{1}{S_i^{z}}{0}\big|^{2},
  \label{eq:Wi-single}
\end{equation}
which is the form used operationally. Summing over the full multiplet
in~\eqref{eq:Wi} is important: because the triplet is three-fold degenerate,
any individual $\ket{1_m}$ returned by the diagonalizer is defined only up to an
arbitrary unitary rotation within the multiplet, whereas the sum in
Eq.~\eqref{eq:Wi} is invariant. Indeed, since $\ket{0}$ is a spin singlet
($S=0,\,S_z=0$), the selection rules make the assignment transparent:
$S_i^{z}$ connects $\ket{0}$ only to the $S_z=0$ triplet member, while
$S_i^{x},S_i^{y}$ (equivalently $S_i^{\pm}=S_i^{x}\pm i S_i^{y}$) connect it to the
$S_z=\pm1$ members. The multiplet sum therefore collects the complete spin-flip
spectral weight localized on site $i$,
\begin{equation}
  W_i \;=\; \sum_{m}\melem{0}{\vec S_i}{1_m}\cdot\melem{1_m}{\vec S_i}{0},
\end{equation}
a real, rotationally invariant quantity that plays the role of a local
``inelastic magnetic susceptibility'' of the transition. For the $N=3$ chain the
$W_i$ are largest on the triangulene sites $T_a,T_b$ (which carry the dominant
$S=1$ character) and are symmetric under the mirror operation of the chain,
consistent with the symmetry of Eq.~\eqref{eq:H}.

\section{Phenomenological spatial mapping (Tersoff--Hamann)}

To turn the site-resolved weights $W_i$ into a real-space image we adopt the
Tersoff--Hamann description of STM, in which the tunneling current at tip
position $\vec r_{\text{tip}}$ is proportional to the local density of states of
the sample evaluated at the tip. For an $s$-wave tip the sample orbitals decay
exponentially into the vacuum barrier,
\begin{equation}
  \psi(\vec r)\;\sim\;\exp\!\big(-\kappa\,|\vec r-\vec r_{\alpha}|\big),
  \qquad
  \kappa=\frac{\sqrt{2m_e\phi}}{\hbar},
  \label{eq:decay}
\end{equation}
where $\phi$ is the effective work function and $\kappa$ the inverse vacuum decay
length (of order $1~\text{\AA}^{-1}$). Each magnetic site $i$ is spatially
distributed over the carbon atoms $\alpha$ that build the corresponding
$\pi$-radical (one phenalenyl, or one half of a triangulene). Placing the tip at
constant height $z_0$ above the molecular plane, $\vec r_{\text{tip}}=(x,y,z_0)$
and $\vec r_\alpha=(x_\alpha,y_\alpha,0)$, and weighting each atomic contribution
by the transition weight of its site, the simulated inelastic conductance map is
\begin{equation}
  \boxed{\;
  \frac{dI}{dV}(x,y)\;\propto\;
    \sum_{i=0}^{4N-1} W_i
    \sum_{\alpha\,\in\,\text{site }i}
    \exp\!\Big(-\kappa\sqrt{(x-x_\alpha)^2+(y-y_\alpha)^2+z_0^2}\,\Big).
  \;}
  \label{eq:map}
\end{equation}
Equation~\eqref{eq:map} is the central phenomenological object: the spin physics
enters entirely through the weights $W_i$ of Eq.~\eqref{eq:Wi}, while the
exponential kernel of Eq.~\eqref{eq:decay} supplies the spatial broadening and the
constant-height blurring set by $z_0$. Larger $\kappa$ or smaller $z_0$ sharpen
the image toward the atomic positions; smaller $\kappa$ or larger $z_0$ produce
the smooth lobes characteristic of experimental constant-height maps.

\paragraph{Bias condition.}
The map~\eqref{eq:map} is the spatial fingerprint of the \emph{single}
excitation of energy $\Delta$. It corresponds to tuning the STM bias to the
inelastic threshold,
\begin{equation}
  e\,V \;=\; \Delta \;=\; E_1-E_0,
\end{equation}
i.e.\ the voltage at which the singlet$\rightarrow$triplet spin-flip channel
opens. Recording $dI/dV$ at this bias and rastering the tip over the molecule
yields, within this model, the image predicted by Eq.~\eqref{eq:map}, directly
comparable to LT-STM IETS data.

\end{document}
